\section{Open problems}\label{sec:conclusion}

\begin{enumerate}
\item \emph{Cryptanalysis of the hidden degree.} The scheme stands or falls
  with the claim that the public key gives no handle on $f$ beyond
  exhaustive search (\S\ref{sec:key}). In particular, the attacker knows the
  action of $\tau$ on $\EA[2^E]$ up to the scalar $f\bmod 2^E$. Independent
  analysis of this point is the most important next step.
\item \emph{Mask entropy.} The binding attack guesses $b/a$ modulo about
  $\sqrt d$. Any way to use fewer torsion bits, for example a dimension-4
  variant of the torsion attack, or a lollipop built from another known
  endomorphism, would lower the security level directly.
\item \emph{Smaller primes.} Decrypting with two half-chains would shrink
  $2^E$ from about $7d^2f^2$ to its square root and cut $p$ from about $d^5$ to
  about $d^3$, by the derivation of \S\ref{sec:params} with $2^{2E}>\deg\gamma$.
\item \emph{Speed.} Theta-coordinate $(2,2)$-chains~\cite{DMPR23} instead of
  Mumford coordinates, stack-only buffers, and SIMD field arithmetic.
\item \emph{Constant-time key generation}, in particular a constant-time
  norm-equation solver.
\item \emph{Compression} of public keys and ciphertexts, by encoding the
  points relative to canonical bases as in compressed SIDH.
\end{enumerate}
